\section{Problem Formulation}
\label{sec:problem}

\paragraph{Setup.}
A generator $G$ produces a document $x = (x_1, \dots, x_n)$ under a watermark key $k$ drawn from a keyspace $\mathcal{K}$. A detector $D$ takes a document and a key and returns a score whose null distribution is known under the unwatermarked hypothesis; we use the $p$-value $p_k(x)$ reported by the same-config test, so that small values indicate the presence of the mark. The detector is valid only against the scheme configuration and key family used at generation, a limitation we return to in \S\ref{sec:limitations}.

\paragraph{Attack catalogue.}
Let $\mathcal{A} = \{t_1, \dots, t_m\}$ be a catalogue of transformations. Each $t$ maps a document to a document, $t: \mathcal{X} \to \mathcal{X}$, and carries an intensity parameter $\theta_t$ that controls how much it changes. Table~\ref{tab:catalogue} lists the catalogue. A transformation may be stochastic, in which case $t(x)$ denotes a draw from its output distribution.

\paragraph{Chains.}
An attack chain is an ordered tuple $T = (t^{(1)}, \dots, t^{(K)})$ with $K \le 3$, applied left to right:
\begin{equation}
  x^{(0)} = x, \quad x^{(j)} = t^{(j)}\!\left(x^{(j-1)};\, \theta^{(j)}\right), \quad x^{(K)} = T(x).
\end{equation}
Order matters. A paraphrase followed by a round trip yields a different document than a round trip followed by a paraphrase, because each step conditions on the output of the previous one. We therefore treat chains as ordered and search over ordered tuples.

\paragraph{Removal.}
For a document $x$ and a key set $K = \{k_1, \dots, k_L\}$ of independent keys, define the document as removed when the mark fails to be detected under every key at level $\alpha$:
\begin{equation}
  \mathrm{removed}_\alpha(x) = \mathbf{1}\!\left[\min_{k \in K} p_k(x) > \alpha \right].
\end{equation}
Two summaries follow. The \emph{removal rate} is the mean of $\mathrm{removed}_\alpha$ over the evaluation set, reported with a Wilson interval. The \emph{corpus operating point} is the true-positive rate at a $1\%$ false-positive rate. Positives are watermarked documents, attacked or not; negatives are unwatermarked documents. The threshold is the highest detector score at which the negative false-positive rate does not exceed $1\%$. We report both, because a per-document threshold and a corpus-level rate answer different questions: the first asks whether this document still carries a mark, the second asks what a detector tuned for a fixed error budget would conclude.

\paragraph{Preservation.}
Let $S(x, T(x))$ be a semantic-preservation score in $[0,1]$, aggregated from the metrics in \S\ref{sec:metrics}. The floor $\tau$ is fixed before the search and reported with every result. A chain that falls below $\tau$ is inadmissible, whatever its removal rate.

\paragraph{Objective.}
The problem we study is
\begin{equation}
  T^\star = \operatorname*{arg\,max}_{T \in \mathcal{T}_{\le 3}} \; \mathbb{E}\left[\mathrm{removed}_\alpha(T(x))\right]
  \quad \text{subject to} \quad
  \mathbb{E}_{x}\left[S(x, T(x))\right] \ge \tau ,
  \label{eq:objective}
\end{equation}
where $\mathcal{T}_{\le 3}$ is the set of ordered chains of length at most three drawn from $\mathcal{A}$. The constraint is enforced on the mean score; for every chain we recommend we also report the fraction of documents that fall below $\tau$. Equation~\ref{eq:objective} makes the trade-off explicit and prevents the degenerate answer, which is to apply the most destructive transformation available.

\paragraph{Why a floor is necessary.}
\citet{sand} show that for any watermark occupying a small share of the space of high-quality responses, a quality-oracle random walk reaches an unmarked output. The construction needs a quality oracle and enough mixing steps, and it degrades quality to the extent that the walk is imperfect. A study that reports removal without a preservation constraint therefore measures the quality of the attacker's oracle as much as the robustness of the scheme. Fixing $\tau$ and reporting the frontier around it keeps the comparison meaningful.
